\documentclass[12pt]{article}
\usepackage[utf8]{inputenc}
\usepackage[T1]{fontenc}
\usepackage{hyperref}
\title{Cajonos Zapotec Language Resources for K-12 Curriculum Development}
\author{Compiled by Claude Sonnet 5.5 (claude-sonnet-5-5) for Andr\'as Kornai}
\date{}
\begin{document}
\maketitle

\section{Introduction}
Cajonos Zapotec (ISO 639-3 zad, Glottocode cajo1238) is an Otomanguean language of Oaxaca, Mexico, with about 5,000 speakers. Searches turned up very little material. No dedicated dictionary, grammar, primer or school textbook was found. Vitality ratings in the sources disagree.

\section{Primary Educational Resources}
No primers, textbooks or school curricula were found in the search results.

\section{Linguistic and Reference Materials}
\begin{itemize}
\item Glottolog language page, which has a bibliography index that can be filtered for dictionaries and grammar sketches. The titles were not visible in the search results: \url{https://glottolog.org/resource/languoid/iso/zad}
\item Hungarian HLT language database entry, which lists two lexical resources and gives about 5,000 first-language speakers: \url{https://hlt.bme.hu/en/dld/language/2989}
\item Glosbe crowdsourced translation pages for Cajonos Zapotec and English. The number of entries was not confirmed: \url{https://glosbe.com/zad/en}
\end{itemize}

\section{Religious and Community Resources}
\begin{itemize}
\item Global Recordings Network page with audio Bible teaching and evangelistic recordings in Cajonos Zapotec, and a link to Scripture resources: \url{https://globalrecordings.net/en/language/zad}
\end{itemize}

\section{Digital and Audio Resources}
The Global Recordings audio above is the only audio found. The Glosbe pages are the only online lexical resource found.

\section{Curriculum Development Context}
The language is thought to be used by the whole community, but it is not known to be taught in schools. A curriculum would need new materials. Local speakers and the sources above would be the starting points. The vitality rating should be checked before planning.
\end{document}